\chapter{Insiemi ordinati: dualità, estremi, diagrammi di Hasse}\label{cap:insiemi-ordinati}

Riprendiamo gli insiemi ordinati del Capitolo~\ref{cap:ordini} con gli strumenti che ora abbiamo: i numeri naturali, l'induzione
e gli insiemi finiti. Vediamo il \textbf{principio di dualità}, che dimezza il lavoro, l'esistenza di elementi minimali negli insiemi
finiti, il legame tra diagrammi di Hasse e isomorfismi, e infine l'\textbf{estremo inferiore} e l'\textbf{estremo superiore}, che ci
porteranno ai reticoli.

% ──────────────────────────────────────────────────────────────
\section{Il principio di dualità}\label{sec:dualita-ordini}

Se $(s, \rho)$ è un insieme ordinato, anche $(s, \bar{\rho})$ lo è (Proposizione~\ref{prop:duale-proprieta}): è il suo \textbf{duale}, e il suo
diagramma di Hasse è quello di $(s, \rho)$ \textbf{capovolto}. Passando al duale i concetti si scambiano a coppie:
\begin{center}
\renewcommand{\arraystretch}{1.2}
\begin{tabular}{c|c|c|c|c|c}
    $\rho$ & minimo & minimale & minorante & estremo inferiore & copre \\ \hline
    $\bar{\rho}$ & massimo & massimale & maggiorante & estremo superiore & è coperto da \\
\end{tabular}
\end{center}

\dfn[enunciato-duale]{Enunciato duale}{
    Sia $e$ un enunciato sugli insiemi ordinati. L'\textbf{enunciato duale} $\bar{e}$ si ottiene da $e$ sostituendo $\rho$ con $\bar{\rho}$,
    cioè scambiando «minimo» con «massimo», «minimale» con «massimale», «minorante» con «maggiorante» e così via.
}

\begin{concetto}
\begin{Teorema}{Principio di dualità per gli insiemi ordinati}{dualita-ordini}
    \begin{prerequisiti}
        \dip{Definizione}{def:enunciato-duale}
        \dip{Proposizione}{prop:duale-proprieta}
    \end{prerequisiti}
    Se un enunciato $e$ vale per tutti gli insiemi ordinati, allora anche il suo duale $\bar{e}$ vale per tutti gli insiemi ordinati.
\end{Teorema}
\begin{dimostrazione}
    Sia $(s, \rho)$ un insieme ordinato. Anche $(s, \bar{\rho})$ è un insieme ordinato, quindi $e$ vale per $(s, \bar{\rho})$. Ma «$e$ per $(s, \bar{\rho})$» è
    proprio «$\bar{e}$ per $(s, \rho)$», perché $\bar{\bar{\rho}} = \rho$: ciò che in $(s, \bar{\rho})$ è un minimo, in $(s, \rho)$ è un massimo, e così via.
    Dunque $\bar{e}$ vale per $(s, \rho)$, che era arbitrario.
\end{dimostrazione}
\end{concetto}

\begin{spiegazione}
    Un teorema vale per \textbf{tutti} gli insiemi ordinati, quindi anche per quelli «capovolti»; e leggere un teorema su un insieme
    capovolto è come leggere il teorema duale sull'insieme di partenza. Per questo, ad esempio, nel Teorema~\ref{teo:unicita-minimo}
    è bastato dimostrare l'unicità del minimo: quella del massimo è l'enunciato duale.
    \begin{center}
    \begin{tikzpicture}[hasse]
        \node[hv] (1) at (0,0) {$1$}; \node[hv] (2) at (-0.8,1.0) {$2$}; \node[hv] (3) at (0.8,1.0) {$3$}; \node[hv] (6) at (0,2.0) {$6$};
        \draw[hlato] (1) -- (2) -- (6); \draw[hlato] (1) -- (3) -- (6);
        \node[font=\small, align=center] at (0,-0.75) {divisori di $6$ con $\mid$\\ minimo $1$, massimo $6$};
        \draw[-{Stealth}, line width=1pt, primary] (1.6,1.0) -- node[above, font=\small] {capovolgo} (3.2,1.0);
        \begin{scope}[xshift=4.8cm]
            \node[hv] (a) at (0,2.0) {$1$}; \node[hv] (b) at (-0.8,1.0) {$2$}; \node[hv] (c) at (0.8,1.0) {$3$}; \node[hv] (d) at (0,0) {$6$};
            \draw[hlato] (a) -- (b) -- (d); \draw[hlato] (a) -- (c) -- (d);
            \node[font=\small, align=center] at (0,-0.75) {stesso insieme con la duale\\ minimo $6$, massimo $1$};
        \end{scope}
    \end{tikzpicture}
    \end{center}
\end{spiegazione}

% ──────────────────────────────────────────────────────────────
\section{Elementi minimali negli insiemi finiti}\label{sec:minimali-finiti}

\begin{concetto}
\begin{Teorema}{Gli insiemi ordinati finiti hanno minimali e massimali}{minimali-esistono}
    \begin{prerequisiti}
        \dip{Definizione}{def:minimo-minimale}
        \dip{Lemma}{lem:togli-elemento}
        \dip{Corollario}{cor:induzione-uso}
        \dip{Teorema}{teo:dualita-ordini}
    \end{prerequisiti}
    Ogni insieme ordinato finito e non vuoto ha almeno un elemento minimale e almeno un elemento massimale.
\end{Teorema}
\begin{dimostrazione}
    Per il principio di dualità basta dimostrare l'esistenza di un minimale. Procediamo per induzione su $n = |s| \geq 1$.
    \textbf{Base}: se $s = \{x\}$, $x$ è minimale. \textbf{Passo}: sia $|s| = n + 1 \geq 2$, sia $x \in s$ e $t = s \setminus \{x\}$, che ha $n \geq 1$
    elementi. Per ipotesi induttiva $t$ (con l'ordine indotto) ha un elemento minimale $m$.
    \begin{caso}{Se $\neg(x \rho m)$}
        Allora $m$ è minimale anche in $s$: l'unico elemento in più rispetto a $t$ è $x$, che non sta sotto $m$.
    \end{caso}
    \begin{caso}{Se $x \rho m$ (e $x \neq m$)}
        Allora $x$ è minimale in $s$: se fosse $y \rho x$ con $y \neq x$, sarebbe $y \in t$ e, per transitività, $y \rho m$ con $y \neq m$
        (se fosse $y = m$, da $m \rho x$ e $x \rho m$ seguirebbe $x = m$), contro la minimalità di $m$ in $t$.
    \end{caso}
\end{dimostrazione}
\end{concetto}

\begin{spiegazione}
    Si costruisce l'insieme un elemento alla volta, tenendo aggiornato un «candidato minimale»: quando arriva un elemento nuovo,
    o non sta sotto il candidato (e il candidato resta minimale), o ci sta (e diventa lui il nuovo candidato). Nell'insieme
    infinito $(\mbZ, \leq)$ il procedimento non finisce mai, e infatti non ci sono minimali: sotto ogni intero ce n'è un altro.
\end{spiegazione}

\begin{concetto}
\begin{Proposizione}{Unico minimale in un insieme finito}{unico-minimale}
    \begin{prerequisiti}
        \dip{Teorema}{teo:minimali-esistono}
        \dip{Proposizione}{prop:ordine-indotto}
        \dip{Lemma}{lem:parti-finite}
    \end{prerequisiti}
    Se un insieme ordinato finito ha un solo elemento minimale, questo è il minimo. Dualmente, un unico elemento massimale di un
    insieme ordinato finito è il massimo.
\end{Proposizione}
\begin{dimostrazione}
    \facoltativa
    Sia $m$ l'unico minimale e sia $x \in s$. L'insieme $t = \{y \in s \mid y \rho x\}$ è finito e non vuoto ($x \in t$), quindi ha un minimale $m'$
    (Teorema~\ref{teo:minimali-esistono}). $m'$ è minimale anche in $s$: se $z \rho m'$, allora $z \rho x$, cioè $z \in t$, e quindi $z = m'$. Per l'unicità
    $m' = m$, e $m = m' \in t$ significa $m \rho x$. Dunque $m$ sta sotto ogni elemento: è il minimo.
\end{dimostrazione}
\end{concetto}

\begin{spiegazione}
    In un insieme finito, scendendo da un elemento qualunque si arriva sempre a un minimale. Se il minimale è uno solo, da ogni
    elemento si scende fino a lui: quindi sta sotto a tutti. Negli insiemi \textbf{infiniti} si può invece scendere (o salire) per sempre.
    Ad esempio, aggiungendo a $(\mbN, \leq)$ un elemento $x$ non confrontabile con nessun numero, $x$ è l'unico massimale (i numeri
    non lo sono: sopra ognuno c'è il successivo), ma non è il massimo:
    \begin{center}
    \begin{tikzpicture}[hasse]
        \foreach \n in {0,...,3} \node[hv] (n\n) at (0,\n*0.75) {$\n$};
        \draw[hlato] (n0) -- (n1) -- (n2) -- (n3); \draw[hlato, dashed] (n3) -- (0,3.4); \node at (0,3.65) {$\vdots$};
        \node[hv evid] (x) at (1.6,1.1) {$x$};
        \node[font=\small, align=left, anchor=west] at (2.4,1.4) {$s = \mbN \cup \{x\}$: la catena dei naturali sale per sempre;\\
            $x$ è isolato: è massimale (e minimale), ma non è il massimo.};
    \end{tikzpicture}
    \end{center}
\end{spiegazione}

% ──────────────────────────────────────────────────────────────
\section{Diagrammi di Hasse e isomorfismi}\label{sec:hasse-isomorfismo}

Negli insiemi ordinati \textbf{finiti} il diagramma di Hasse contiene tutta l'informazione sull'ordine. La chiave è il lemma seguente.

\begin{concetto}
\begin{Lemma}{Catene di coperture}{catene-copertura}
    \begin{prerequisiti}
        \dip{Definizione}{def:copertura}
        \dip{Teorema}{teo:ordine-finito}
        \dip{Teorema}{teo:induzione-seconda}
    \end{prerequisiti}
    Sia $(s, \leq)$ un insieme ordinato finito e siano $x < y$. Allora esistono $c_0 = x, c_1, \dots, c_k = y$ ($k \geq 1$) tali che $c_{i + 1}$ copre $c_i$
    per ogni $i$: da $x$ si arriva a $y$ salendo lungo i segmenti del diagramma di Hasse.
\end{Lemma}
\begin{dimostrazione}
    Per induzione di seconda forma sul numero $n$ di elementi dell'insieme $(x, y) := \{z \in s \mid x < z < y\}$. Se $n = 0$, $y$ copre $x$: basta $k = 1$.
    Se $n > 0$, sia $z \in (x, y)$. Gli insiemi $(x, z)$ e $(z, y)$ sono parti proprie di $(x, y)$ (non contengono $z$), quindi hanno meno di $n$
    elementi (Teorema~\ref{teo:ordine-finito}). Per ipotesi induttiva c'è una catena di coperture da $x$ a $z$ e una da $z$ a $y$: unendole si ottiene
    una catena da $x$ a $y$.
\end{dimostrazione}
\end{concetto}

\begin{spiegazione}
    Tra $x$ e $y$ ci sono finitamente molti elementi; se non ce n'è nessuno, $y$ copre $x$; altrimenti se ne sceglie uno in mezzo e si
    ripete il ragionamento sui due tratti, più corti. Nei divisori di $12$: da $1$ a $12$ si sale per $1 \to 2 \to 4 \to 12$ oppure $1 \to 3 \to 6 \to 12$.
    In $(\mbQ, \leq)$ invece il lemma è falso: nessun razionale copre un altro (tra due razionali ce n'è sempre un terzo), quindi il diagramma
    di Hasse di $(\mbQ, \leq)$ non ha nessun segmento.
\end{spiegazione}

\begin{concetto}
\begin{Teorema}{Isomorfismi e diagrammi di Hasse}{hasse-isomorfismo}
    \begin{prerequisiti}
        \dip{Lemma}{lem:catene-copertura}
        \dip{Definizione}{def:crescente}
        \dip{Definizione}{def:copertura}
    \end{prerequisiti}
    Siano $(s, \rho)$ e $(s', \rho')$ insiemi ordinati finiti e $f \colon s \to s'$ biettiva. Allora $f$ è un isomorfismo se e solo se rispetta la
    relazione di copertura: per ogni $x, y \in s$, $y$ copre $x$ se e solo se $f(y)$ copre $f(x)$. In altre parole, due insiemi ordinati finiti
    sono isomorfi se e solo se hanno lo stesso diagramma di Hasse.
\end{Teorema}
\begin{dimostrazione}
    \begin{caso}{$(\rightarrow)$ Un isomorfismo rispetta le coperture}
        «$y$ copre $x$» significa: $x \neq y$, $x \rho y$ e non esiste $z \neq x, y$ con $x \rho z \rho y$. Poiché $f$ è biettiva e $x \rho y \leftrightarrow f(x) \rho' f(y)$, ogni
        ingrediente si trasporta: $x \neq y \leftrightarrow f(x) \neq f(y)$; un $z$ tra $x$ e $y$ dà $f(z)$ tra $f(x)$ e $f(y)$, e viceversa un $z'$ tra $f(x)$ e $f(y)$
        è $f(z)$ per un $z$ (suriettività) che sta tra $x$ e $y$. Quindi $y$ copre $x$ se e solo se $f(y)$ copre $f(x)$.
    \end{caso}
    \begin{caso}{$(\leftarrow)$ Rispettare le coperture basta}
        Per il Lemma~\ref{lem:catene-copertura}, $x \rho y$ se e solo se $x = y$ oppure c'è una catena di coperture da $x$ a $y$. Poiché $f$ è biettiva e
        porta catene di coperture in catene di coperture (e $f^{-1}$ fa il viceversa), $x \rho y$ se e solo se $f(x) = f(y)$ oppure c'è una catena di
        coperture da $f(x)$ a $f(y)$, cioè se e solo se $f(x) \rho' f(y)$.
    \end{caso}
\end{dimostrazione}
\end{concetto}

\begin{spiegazione}
    Per decidere se due insiemi ordinati finiti sono isomorfi basta disegnarne i diagrammi di Hasse e vedere se sono lo stesso disegno a
    meno dei nomi. Ad esempio i divisori di $30$ con $\mid$ e le parti di $\{a, b, c\}$ con $\subseteq$ hanno lo stesso diagramma (un cubo):
    \begin{center}
    \begin{tikzpicture}[hasse]
        \node[hv] (1) at (0,0) {$1$};
        \node[hv] (2) at (-1.2,1.1) {$2$}; \node[hv] (3) at (0,1.1) {$3$}; \node[hv] (5) at (1.2,1.1) {$5$};
        \node[hv] (6) at (-1.2,2.2) {$6$}; \node[hv] (10) at (0,2.2) {$10$}; \node[hv] (15) at (1.2,2.2) {$15$};
        \node[hv] (30) at (0,3.3) {$30$};
        \draw[hlato] (1) -- (2); \draw[hlato] (1) -- (3); \draw[hlato] (1) -- (5);
        \draw[hlato] (2) -- (6); \draw[hlato] (2) -- (10); \draw[hlato] (3) -- (6); \draw[hlato] (3) -- (15); \draw[hlato] (5) -- (10); \draw[hlato] (5) -- (15);
        \draw[hlato] (6) -- (30); \draw[hlato] (10) -- (30); \draw[hlato] (15) -- (30);
        \draw[-{Stealth}, line width=1pt, primary] (2.1,1.65) -- node[above, font=\small, align=center] {$2 \mapsto \{a\}$, $3 \mapsto \{b\}$,\\ $5 \mapsto \{c\}$} (5.5,1.65);
        \begin{scope}[xshift=7.6cm]
            \node[hv] (v) at (0,0) {$\varnothing$};
            \node[hv] (a) at (-1.5,1.1) {$\{a\}$}; \node[hv] (b) at (0,1.1) {$\{b\}$}; \node[hv] (c) at (1.5,1.1) {$\{c\}$};
            \node[hv] (ab) at (-1.5,2.2) {$\{a, b\}$}; \node[hv] (ac) at (0,2.2) {$\{a, c\}$}; \node[hv] (bc) at (1.5,2.2) {$\{b, c\}$};
            \node[hv] (abc) at (0,3.3) {$\{a, b, c\}$};
            \draw[hlato] (v) -- (a); \draw[hlato] (v) -- (b); \draw[hlato] (v) -- (c);
            \draw[hlato] (a) -- (ab); \draw[hlato] (a) -- (ac); \draw[hlato] (b) -- (ab); \draw[hlato] (b) -- (bc); \draw[hlato] (c) -- (ac); \draw[hlato] (c) -- (bc);
            \draw[hlato] (ab) -- (abc); \draw[hlato] (ac) -- (abc); \draw[hlato] (bc) -- (abc);
        \end{scope}
    \end{tikzpicture}
    \end{center}
    L'isomorfismo manda ogni divisore nell'insieme dei suoi fattori primi (scritti come $a, b, c$). Per gli insiemi infiniti non basta:
    $(\mbQ, \leq)$ e $\mbQ$ con l'uguaglianza hanno entrambi un diagramma di Hasse senza segmenti, ma il primo è totalmente ordinato e il
    secondo no.
\end{spiegazione}

% ──────────────────────────────────────────────────────────────
\newpage
\section{Ordine indotto da una funzione}\label{sec:ordine-funzione}

Un ordine su un insieme si può «trasportare» su un altro insieme tramite una funzione iniettiva.

\begin{concetto}
\begin{Proposizione}{Ordine indotto da una funzione}{ordine-funzione}
    \begin{prerequisiti}
        \dip{Definizione}{def:tipi-relazioni}
        \dip{Definizione}{def:crescente}
        \dip{Proposizione}{prop:ridotta}
    \end{prerequisiti}
    Siano $a$ un insieme, $(b, \sigma)$ un insieme ordinato e $f \colon a \to b$. Definiamo su $a$ la relazione $\rho_f$ ponendo
    $x \rho_f y \leftrightarrow f(x) \sigma f(y)$. Allora $\rho_f$ è riflessiva e transitiva; è un ordine se e solo se $f$ è iniettiva, e in tal caso la ridotta
    $a \to \mathrm{Im} f$ è un isomorfismo tra $(a, \rho_f)$ e $(\mathrm{Im} f, \sigma)$.
\end{Proposizione}
\begin{dimostrazione}
    $f(x) \sigma f(x)$ dà la riflessività; da $f(x) \sigma f(y)$ e $f(y) \sigma f(z)$ segue $f(x) \sigma f(z)$: transitività.
    \begin{caso}{Se $f$ è iniettiva, $\rho_f$ è antisimmetrica}
        Da $x \rho_f y$ e $y \rho_f x$ segue $f(x) \sigma f(y)$ e $f(y) \sigma f(x)$, quindi $f(x) = f(y)$ e $x = y$.
    \end{caso}
    \begin{caso}{Se $f$ non è iniettiva, $\rho_f$ non è antisimmetrica}
        Se $f(x) = f(y)$ con $x \neq y$, allora $x \rho_f y$ e $y \rho_f x$.
    \end{caso}
    Infine, se $f$ è iniettiva la ridotta è biettiva e, per definizione, $x \rho_f y \leftrightarrow f(x) \sigma f(y)$: è un isomorfismo.
\end{dimostrazione}
\end{concetto}

\begin{spiegazione}
    Si ordinano gli elementi di $a$ \textbf{confrontando le loro «etichette»} $f(x)$. Ordinare gli studenti per numero di matricola
    funziona (matricole diverse per studenti diversi: $f$ è iniettiva). Ordinare le parole per lunghezza no: «sole» e «mare» hanno la stessa
    lunghezza, quindi ciascuna starebbe sotto l'altra pur essendo diverse, e l'antisimmetria cade.
\end{spiegazione}

% ──────────────────────────────────────────────────────────────
\section{Estremo inferiore ed estremo superiore}\label{sec:inf-sup}

\dfn[inf-sup]{Estremo inferiore ed estremo superiore}{
    Sia $(s, \rho)$ un insieme ordinato e sia $t \subseteq s$. Se l'insieme $\mathrm{Minor}(t)$ dei minoranti di $t$ ha massimo, questo si chiama
    \textbf{estremo inferiore} di $t$ e si scrive $\inf(t)$ (o $\inf_{(s, \rho)}(t)$). Dualmente, se $\mathrm{Maggior}(t)$ ha minimo, questo si chiama
    \textbf{estremo superiore} di $t$ e si scrive $\sup(t)$.
}

L'estremo inferiore è il \textbf{più grande} degli elementi che stanno sotto a tutto $t$; non è detto che stia in $t$. Se ci sta, è il minimo.

\begin{concetto}
\begin{Proposizione}{Minimo ed estremo inferiore}{minimo-inf}
    \begin{prerequisiti}
        \dip{Definizione}{def:inf-sup}
        \dip{Definizione}{def:minoranti}
    \end{prerequisiti}
    Se $t$ ha minimo $m$, allora $\inf(t) = m$. Dualmente, se $t$ ha massimo $M$, allora $\sup(t) = M$. In particolare $\inf(\{a\}) = \sup(\{a\}) = a$.
\end{Proposizione}
\begin{dimostrazione}
    $m$ sta sotto ogni elemento di $t$, quindi è un minorante. Ogni minorante $z$ sta sotto ogni elemento di $t$, in particolare sotto
    $m \in t$: quindi $m$ è il massimo dei minoranti.
\end{dimostrazione}
\end{concetto}

\begin{spiegazione}
    Sul diagramma dei divisori di $12$, per $t = \{4, 6\}$: i minoranti sono $1$ e $2$, e il più alto è $2$, quindi $\inf\{4, 6\} = 2$ (il massimo
    comun divisore); i maggioranti sono solo $12$, quindi $\sup\{4, 6\} = 12$ (il minimo comune multiplo). Per $t = \{2, 4\}$, che ha minimo $2$
    e massimo $4$: $\inf = 2$, $\sup = 4$.
    \begin{center}
    \begin{tikzpicture}[hasse]
        \node[hv] (1) at (0,0) {$1$}; \node[hv evid2] (2) at (-0.8,1.0) {$2$}; \node[hv] (3) at (0.8,1.0) {$3$};
        \node[hv evid] (4) at (-1.6,2.0) {$4$}; \node[hv evid] (6) at (0,2.0) {$6$}; \node[hv evid2] (12) at (-0.8,3.0) {$12$};
        \draw[hlato] (1) -- (2) -- (4) -- (12); \draw[hlato] (1) -- (3) -- (6) -- (12); \draw[hlato] (2) -- (6);
        \node[font=\small, align=left, anchor=west] at (1.6,1.5) {in arancione $t = \{4, 6\}$;\\ in verde $\inf(t) = 2$ e $\sup(t) = 12$.};
    \end{tikzpicture}
    \end{center}
\end{spiegazione}

\begin{esempio}{Estremi in insiemi ordinati notevoli}
    \begin{itemize}
        \item In $(\mbN, \leq)$ ogni parte non vuota $t$ ha estremo inferiore $\min(t)$; ha estremo superiore se e solo se è superiormente
              limitata, e in tal caso $\sup(t) = \max(t)$ (ordine naturale).
        \item In $(\mbN, \mid)$ l'estremo inferiore di $\{a, b\}$ è il massimo comun divisore non negativo e l'estremo superiore il minimo comune
              multiplo non negativo (Proposizione~\ref{prop:mcd-associati}): ad esempio $\inf\{12, 18\} = 6$ e $\sup\{12, 18\} = 36$.
        \item In $(\mcP(s), \subseteq)$, per ogni parte non vuota $t$ di $\mcP(s)$: $\inf(t) = \bigcap t$ e $\sup(t) = \bigcup t$
              (Proposizione~\ref{prop:inf-parti}).
    \end{itemize}
\end{esempio}

\begin{concetto}
\begin{Proposizione}{Estremi nell'insieme delle parti}{inf-parti}
    \begin{prerequisiti}
        \dip{Definizione}{def:inf-sup}
        \dip{Definizione}{def:intersezione}
        \dip{Proposizione}{prop:appartenenza}
    \end{prerequisiti}
    Sia $s$ un insieme e sia $t$ una parte non vuota di $\mcP(s)$. In $(\mcP(s), \subseteq)$ si ha $\inf(t) = \bigcap t$ e $\sup(t) = \bigcup t$.
\end{Proposizione}
\begin{dimostrazione}
    \begin{caso}{$\inf(t) = \bigcap t$}
        $\bigcap t$ è contenuto in ogni elemento di $t$, quindi è un minorante. Se $z \subseteq s$ è un minorante, cioè $z \subseteq y$ per ogni $y \in t$,
        ogni elemento di $z$ sta in tutti gli $y \in t$, cioè in $\bigcap t$: $z \subseteq \bigcap t$. Quindi $\bigcap t$ è il massimo dei minoranti.
    \end{caso}
    \begin{caso}{$\sup(t) = \bigcup t$}
        $\bigcup t \subseteq s$ contiene ogni elemento di $t$; se $z$ contiene ogni elemento di $t$, contiene ogni elemento di ogni $y \in t$, cioè
        $\bigcup t \subseteq z$. Quindi $\bigcup t$ è il minimo dei maggioranti.
    \end{caso}
\end{dimostrazione}
\end{concetto}

\begin{spiegazione}
    Il più grande insieme contenuto in tutti gli elementi di $t$ è la loro parte comune; il più piccolo insieme che li contiene tutti è la
    loro unione. Con $t = \{\{1, 2\}, \{2, 3\}\}$: $\inf(t) = \{2\}$ e $\sup(t) = \{1, 2, 3\}$.
\end{spiegazione}
